\section{Numerical Implementation}
\label{sec:numerics}
\label{sec:implementation}

The model is implemented as a modular simulation engine in C (ISO~C99) and
built with CMake. Environmental forcing, gas exchange, carbohydrate dynamics,
nutrient uptake and biomass growth are written as separate translation units.
Each physiological block can therefore be tested on its own, while the coupled
path from forcing to growth stays intact. This section describes how the state
equations are advanced, and which checks were applied to the solver components
before any result was produced.

\subsection{Simulation architecture and time integration}

Each right-hand-side evaluation follows the fast--slow structure of
Section~\ref{sec:model} and Figure~\ref{fig:overview}. Four blocks are executed
in sequence. The first block
evaluates the clock time and the environmental drivers: PAR, leaf temperature,
relative humidity, vapour-pressure deficit, ambient CO$_2$ and leaf water
potential. The second block computes the gas-exchange layer: stomatal
conductance, the electron transport rate and the two Farquhar branches. The
intercellular CO$_2$ concentration follows from the closed-form solution of
Eq.~\eqref{eq:ci_quadratic}. The third block converts net assimilation into the
carbon-input term $\Phi_c$. The fourth block assembles the derivatives of the
eleven slow states. These are starch, cytosolic and vacuolar sucrose, the
partition and allocation variables, the two nutrient pools, the two affinity
states, and leaf and root biomass.

Gas exchange is algebraic and is solved at every stage of every step, not once
per output sample. This matters for the coupling. Conductance, assimilation and
$C_i$ are functions of the current state and the current forcing, so a stage
value taken from a stale operating point would break the supply--demand balance
that the closure enforces. The eleven slow states carry no algebraic
constraint and are integrated as an explicit system.

Integration uses a separate, reusable C ODE-solver library through a single
interface, so the model does not depend on one backend. The production
configuration selects an adaptive Dormand--Prince embedded pair
(\texttt{ODE45}) with an error tolerance of $10^{-8}$
\citep{HairerNorsettWanner1993}. Under an adaptive method the configuration
entry \texttt{dt\_hours} is not the integration step. It sets the initial trial
step, which the step controller immediately overrides, and it sets the interval
at which the trajectory is sampled for output. Production runs use
$\texttt{dt\_hours}=0.1$~h, so the reported trajectory is a fixed
0.1~h grid over a 30-day horizon, while the accuracy of that trajectory is
controlled by the tolerance alone.

A fixed step is used where the adaptive controller is either unusable or
unnecessary. The sensitivity and sweep runners override the backend to
fixed-step RK4, because those designs draw parameter values at the edges of the
declared ranges, where the adaptive step can collapse. A fixed step also gives a
predictable cost over several thousand evaluations. The carbon-balance
tool of Section~\ref{sec:carbon_closure} uses fixed-step RK4 for a different
reason: only a fixed-step driver can advance the flux accumulators with the same
weights and the same step as the states. Both backends are run through the same
production path at the same settings, so the effect of the choice can be
measured rather than assumed.

Two runtime guards protect batch work. A non-finite guard and a domain guard
observe the trajectory and classify a run as divergent. A step-budget plugin
counts attempted steps and stops the run when a generous budget is exhausted.
The second guard covers a failure the first cannot see: for some parameter
combinations the adaptive step collapses while the state stays finite, so the
integrator neither advances nor reports an error. Such a run then unwinds
normally and is recorded as inadmissible, which is how the batch drivers already
treat divergence.

\subsection{Verification of the solver components}
\label{sec:verification}

The gas-exchange closure was verified by substitution rather than by assertion.
For a reference leaf and PAR from 10 to
2000~$\mu$mol~m$^{-2}$~s$^{-1}$, the analytic solution was inserted back into
the supply relation and the residual $|A_n - g_c(C_a - C_i)|$ was recorded. The
analytic branch solution reproduces the intended operating point. The residual
stays at machine precision over the whole light range, so the closed-form root
solves the coupled system rather than approximating it.

Two limits were checked separately, because both are boundaries where an
implementation can silently disagree with the equations. When conductance
reaches zero the implementation returns $C_i = \Gamma^\ast$ and
$A_n = -R_d$, which is the dark limit of the formulation. Above that limit,
assimilation crosses zero smoothly at $\mathrm{PAR}=13.36$~$\mu$mol~m$^{-2}$~s$^{-1}$
under the reference conditions. An earlier hard cut-off at
$\mathrm{PAR}=20$~$\mu$mol~m$^{-2}$~s$^{-1}$ was removed, because $f_{PAR}$ is
not zero there and the returned assimilation jumped across the threshold. With
the analytic closure the terminal $\max(0,\cdot)$ clamp already returns exactly
zero below the light compensation point, continuously and without a special
case. Removing the gate therefore also removes a discontinuity that an adaptive
step controller has to resolve twice per simulated day.

\subsection{Carbohydrate-subsystem bookkeeping closure}
\label{sec:carbon_closure}

Whether the carbon bookkeeping closes cannot be answered by post-processing the
solver output. The trajectory records pool sizes, whereas closure is a statement
about time-integrated fluxes, and a quadrature applied afterwards would test the
quadrature as much as the model. The carbon-balance tool therefore appends ten
flux accumulators to the eleven states and advances the augmented system with
identical RK4 weights and an identical step. Both parts come from one call to the
right-hand side, so the accumulators see the same stage values as the state
integration. The resulting budget is then exact to the arithmetic of the
integrator.

The closure is asserted for the carbohydrate subsystem
$\{a, s, s_v\}$, that is starch, cytosolic sucrose and vacuolar sucrose. Adding
Eqs.~\eqref{eq:starch}, \eqref{eq:sucrose} and \eqref{eq:sucrose_vac} makes the
partition terms $k_s \Phi_c$ and $(1-k_s)\Phi_c$ recombine, the degradation flux
$D_a$ cancel between the starch and sucrose equations, and the vacuolar exchange
$k_{\mathrm{in}} s - k_{\mathrm{out}} s_v$ cancel between the sucrose and
vacuolar equations. Integrating what remains over $[0,T]$ gives
\begin{equation}
\bigl(a + s + s_v\bigr)\Big|_0^{T}
=
\int_0^{T} \Phi_c \, dt
-
\int_0^{T}
\bigl( C_u + C_t + R_m + \eta\,E + C_g \bigr) dt ,
\label{eq:budget}
\end{equation}
with every symbol as defined in Eqs.~\eqref{eq:uptake_cost}--\eqref{eq:growth_cost}
and~\eqref{eq:sucrose}, so
the identity can be read off the model equations without translation. Two points
about the terms. The integrand on the input side is the model carbon input
$\Phi_c$ of Eq.~\eqref{eq:conversion}, in $\mu$mol hexose gFW$^{-1}$ h$^{-1}$,
and not the leaf-level $A_n$: the accumulator integrates the same converted
quantity that drives the pools, so both sides carry one set of units. And the construction
cost $C_g$ is identically zero in every reported run because $Y = 0$, so the
implementation carries no accumulator for it: the ten accumulators are
assimilation, the two partition fluxes, starch degradation, the uptake,
transport, maintenance and loading costs, and the two vacuolar exchange fluxes.
The term is written in the identity because it is required whenever $Y$ is
activated, and a reader checking the residual against the code should expect it
to be absent there for that reason. The residual is the difference between the two sides, reported
relative to cumulative carbon input,
\begin{equation}
\varepsilon_{C}
=
\frac{1}{\int_0^{T} \Phi_c \, dt}
\left[
\bigl(a + s + s_v\bigr)\Big|_0^{T}
-
\int_0^{T} \Phi_c \, dt
+
\int_0^{T}
\bigl( C_u + C_t + R_m + \eta\,E + C_g \bigr) dt
\right].
\label{eq:residual}
\end{equation}
Across the four calibration photoperiods the relative residual lies between
$2.9\times10^{-14}$ and $6.1\times10^{-14}$, which is machine precision for
double-precision accumulation over this many steps. No carbon is created or
destroyed numerically inside the declared scope.

The scope must be stated exactly, because it does not cover the whole plant.
The biomass gain term sits outside the budget by construction. In the growth law
of Eq.~\eqref{eq:dBl} the sucrose drain is
$\phi_{\text{load}} E_{\text{night}}$ while the biomass gain is
$\lambda_{sb}\,\phi_{\text{load}} E_{\text{night}} B_l$, so the two differ by
the factor $\lambda_{sb} B_l$. That factor equals unity only at
$B_l = 1/\lambda_{sb}$, which is of order $10^{2}$~gFW. Over the simulated
seedling range the factor stays $10^{2}$ to $10^{3}$ away from unity, so the two
terms are not two views of one carbon flow. The growth
term is a rate law calibrated against observed RGR, not a stoichiometric
conversion of sucrose into biomass carbon. Eq.~\eqref{eq:residual} therefore
verifies bookkeeping closure of the carbohydrate subsystem, and it is reported
under that name only. A budget that included biomass would require a
stoichiometric growth equation, which the model does not have.

\subsection{Convergence and tolerance study}

Three numerical settings can in principle contaminate a reported number: the
solver tolerance, the output interval and the gas-exchange tolerance. Each was
refined independently around its production value, together with a step sweep of
the fixed-step backend and a comparison of the two backends at production
settings. Every step used divides both 24~h and the photoperiod exactly, so the
end-of-day sample lands on a grid point and no interpolation error enters.

The reported quantities are insensitive to these settings over the tested range.
At the production step of $0.1$~h the worst-case relative deviation from the
$dt\to0$ reference over the four photoperiods is $4.3\times10^{-4}$ in RGR,
$3.1\times10^{-3}$ in final biomass and $2.0\times10^{-3}$ in end-of-day starch;
the cytosolic sucrose pool is the least well resolved at $6.8\times10^{-2}$, as
the smallest and fastest-turning state. The algebraic tolerance inside the $C_i$
solve moves RGR by at most $3.2\times10^{-7}$, and the production configuration
uses the closed-form solver, which has no such tolerance. The measured order of
the fixed-step sweep is close to one rather than four, because starch degradation
is multiplied by the light state and the right-hand side is discontinuous twice
per simulated day. None of these deviations approaches the $10^{-2}$ to
$10^{-1}$ discrepancies the model shows against measurement, so numerical error
is not what limits agreement. The full study, including the Richardson estimates,
the per-photoperiod tables, the legacy iterative comparator and the one external
integrator excluded after it failed a consistency check, is in
\path{docs/revision/convergence_study.md}.

\section{Model Evaluation Protocol}
\label{sec:protocol}

This section defines the evaluation rules and reports no results. The reference
data are stated with their source, the fitted coefficients are counted in each
regime, and every metric is defined together with what it does not measure. No
number in Section~\ref{sec:results} should then be readable as more than it is.

\subsection{Reference data and growth conditions}

Four photoperiods are used for calibration and cross-validation: 4, 6, 8 and
12~h of light per 24~h cycle. An 18~h photoperiod is held out of every fit
entirely and is used only as an out-of-sample condition. The four short
photoperiods span the range over which the reference experiments report growth,
and 18~h lies outside it, which is why it is treated separately.

All reported runs reproduce the growth conditions of the reference experiments.
PAR is 160~$\mu$mol~m$^{-2}$~s$^{-1}$ delivered as a square wave over the light
phase, and leaf temperature is 21~$^\circ$C by day and 19~$^\circ$C by night.
This is a methodological requirement, not a detail. The shipped configuration
instead runs a half-sine profile peaking at
1200~$\mu$mol~m$^{-2}$~s$^{-1}$ at a single temperature. A half-sine averages
$2/\pi$ of its peak, so that daily light integral is about 4.8 times the chamber
value at the same day length. A coefficient fitted under one light dose does not
transfer to another:
the carbon-input coefficient would absorb the dose mismatch, and the resulting
value would describe the forcing rather than the plant. Scoring the model
against chamber data under the shipped forcing would therefore compare a model
driven at one photon dose with plants grown at another.

Every reference value used in this paper comes from one study
\citep{Sulpice2014}, and we state how that was established because an earlier
draft of this work attributed part of the data to the wrong source. The RGR
targets are 0.068, 0.114, 0.171 and 0.260~d$^{-1}$ for the 4, 6, 8 and 12~h
photoperiods. The diel carbohydrate reference is
\path{data/paper/sugars_starch_reference.csv}, which holds 65 rows over five
photoperiods, that is 13 sampling times per photoperiod at a 2~h resolution
through the 24~h cycle, with the row at ZT${}=24$ a copy of ZT${}=0$. Starch and
sucrose are used; the remaining assayed metabolites are not part of this
evaluation.

The diel series were traced to Supplemental Table S4 of that study, the sheet of
replicate means, and independently verified against the published supplement by a
comparison script released with the code, so the identification can be repeated
rather than taken on trust.

That the growth rates, the dusk starch values and the diel series all come from
one growth experiment is what makes the reference set internally consistent: a
model scored against it is not being asked to reconcile two laboratories. Two
other studies of diel carbohydrate behaviour cover only a subset of these day
lengths, at 12~h \citep{Annunziata2017} and at 6 and 12~h \citep{Mengin2017}, and
neither can supply the four-point series.



\subsection{Parameterisation regimes}

Three parameterisation regimes are evaluated (Table~\ref{tab:scenarios}). They
differ only in how many coefficients are fitted and in whether those
coefficients are allowed to depend on photoperiod. Everything else, including
the forcing, the initial state and all sourced parameters, is identical.

The saturated regime is the legacy per-photoperiod baseline. It fits two coefficients per
photoperiod by nested bisection: the growth coefficient $\lambda_{sb}$ against
the observed RGR, and the daytime starch synthesis rate $\lambda_{sdi}$ against
the observed end-of-day starch. With four photoperiods that is eight fitted
values against eight observations, so the regime has zero degrees of freedom.
Its agreement with the data cannot be reported as predictive accuracy. What a
bisection leaves behind is a convergence residual of the search, controlled by
the resolution of the parameter grid, and it goes to zero as the search is
tightened. The regime is retained so that the published result stays
reproducible, and it is reported as a calibration diagnostic only.

The shared regime fits four coefficients for all photoperiods together:
$\lambda_{sb}$, $\lambda_{sdr}$, $\lambda_{sdi}$ and the sink activation
constant. Photoperiod is then the only thing that differs between the four
runs, so any photoperiod response the model produces is generated by its
dynamics rather than supplied by its parameters. The third regime keeps the
same shared set but allows the growth coefficient to vary linearly with
photoperiod $P$,
\begin{equation}
\lambda_{sb}(P) = c_0 + c_1 P ,
\label{eq:lsb_linear}
\end{equation}
which replaces one shared constant by two coefficients. Eq.~\eqref{eq:lsb_linear}
encodes the hypothesis that the growth response to photoperiod is smooth in $P$
and is not fully generated by the coupled dynamics. It remains a prediction for
a held-out photoperiod, because the held-out observation never enters the fit.

\begin{table}[!htbp]
\centering
\small
\caption{Parameterisation regimes. All three use the same chamber forcing
(PAR 160~$\mu$mol~m$^{-2}$~s$^{-1}$, square wave, 21/19~$^\circ$C), the same
initial state and the same sourced parameters. $P$ denotes photoperiod in
hours. The observation count is two per photoperiod, the published RGR and the
published end-of-day starch. Regime R1 is reported as a calibration diagnostic;
only R2 and R3 enter the cross-validation of
Section~\ref{sec:crossval_design}.}
\label{tab:scenarios}
\begin{tabular}{@{}llrrl@{}}
\toprule
\textbf{Regime} & \textbf{Fitted coefficients} & \textbf{Fitted} &
\textbf{Obs.} & \textbf{Role} \\
 & & \multicolumn{1}{c}{\footnotesize values} &
\multicolumn{1}{c}{\footnotesize used} & \\
\midrule
R1 saturated & $\lambda_{sb}$, $\lambda_{sdi}$, per photoperiod & 8 & 8 &
calibration diagnostic \\
R2 shared & $\lambda_{sb}$, $\lambda_{sdr}$, $\lambda_{sdi}$, sink constant & 4
& 8 & held-out prediction \\
R3 shared, linear & as R2 with $\lambda_{sb}=c_0+c_1 P$ & 5 & 8 &
held-out prediction \\
\bottomrule
\end{tabular}
\end{table}

\subsection{Metrics}

Growth is evaluated on relative growth rate, computed from total biomass over
the simulated interval,
\begin{equation}
\mathrm{RGR} = \frac{\ln B(t_2) - \ln B(t_1)}{t_2 - t_1} .
\label{eq:rgr}
\end{equation}
Per photoperiod we report the signed relative error. Over a set of photoperiods
we report the mean absolute percentage error:
\begin{equation}
E_{\text{rel}}
= 100\,\frac{\widehat{\mathrm{RGR}} - \mathrm{RGR}_{\text{ref}}}
{\mathrm{RGR}_{\text{ref}}},
\qquad
\mathrm{MAPE} = \frac{1}{m}\sum_{j=1}^{m} \bigl| E_{\text{rel},j} \bigr| ,
\label{eq:rgr_error}
\end{equation}
where $\widehat{\mathrm{RGR}}$ is the simulated value and $m$ is the number of
photoperiods in the set. In-sample and held-out MAPE are always reported as
separate quantities and never pooled.

The diel carbohydrate series are evaluated with a normalised root-mean-square
error. A single multiplicative scale $\alpha$ is first fitted to each simulated
trace by least squares,
\begin{equation}
\alpha = \frac{\sum_{i=1}^{n} \hat{y}_i y_i}{\sum_{i=1}^{n} \hat{y}_i^{2}},
\qquad
r_i = \alpha \hat{y}_i - y_i ,
\label{eq:alpha}
\end{equation}
and the metrics are then computed on the scaled residual $r_i$,
\begin{equation}
\mathrm{nRMSE} = \frac{1}{|\bar{y}|}
\sqrt{\frac{1}{n}\sum_{i=1}^{n} r_i^{2}},
\qquad
\mathrm{MAE} = \frac{1}{n}\sum_{i=1}^{n} |r_i| .
\label{eq:nrmse}
\end{equation}
Here $y_i$ is the reference value and $\hat{y}_i$ the simulated value, $n$ is
the number of matched sampling times, and $\bar{y}$ is the mean reference
level. MAE
keeps the units of the reference data and is computed for the same fitted
$\alpha$, so both metrics describe one and the same fit.

The consequence of fitting $\alpha$ must be stated, because it bounds what the
carbohydrate comparison can support. A fitted multiplicative scale removes the
absolute level from the residual, so nRMSE scores the shape and the timing of
the diel cycle and not the absolute pool concentration. A trace with the
correct rhythm and a systematically wrong amplitude can pass; a trace with the
correct amplitude and the wrong phase cannot. For that reason $\alpha$ is
reported beside every nRMSE in Section~\ref{sec:results}: a value near unity
means shape and level agree, and a value far from unity localises the
disagreement in the level alone. The acceptance thresholds are
predefined heuristic criteria rather than external standards: $\mathrm{nRMSE}
\le 0.25$ for starch and $\mathrm{nRMSE} \le 0.50$ for sucrose. The sucrose
threshold is the looser of the two because whole-tissue sucrose pools cytosolic
and vacuolar sucrose, and the split between the two compartments is itself
uncertain.

\subsection{Cross-validation design}
\label{sec:crossval_design}

Growth prediction is evaluated by leave-one-photoperiod-out cross-validation.
The split is by photoperiod rather than at random, because interpolation across
the photoperiod range is exactly the ability in question. Four folds are run in
each of the two shared regimes, which gives eight held-out predictions in total.

One coefficient is refitted per fold, and stating this precisely bounds what the
experiment can support. The growth coefficient $\lambda_{sb}$ is searched on a
fixed grid of eleven values from $0.00200$ to $0.00480$~gFW ($\mu$mol
C$_6$)$^{-1}$. For each grid value the model is run at all four
photoperiods, and the value minimising the mean absolute percentage error in RGR
over the three training photoperiods is selected, frozen, and used to predict the
fourth. In regime A the selected value is applied unchanged at the held-out
photoperiod. In regime B the straight line of Eq.~\eqref{eq:lsb_linear}
is fitted to the three per-photoperiod grid optima of the training set by least
squares, and evaluated at the held-out photoperiod. Regime B therefore carries two coefficients where regime A carries
one.

The other three calibrated coefficients, $\lambda_{sdr}$, $\lambda_{sdi}$ and
$K_s$, are not refitted in any fold: they keep the values obtained from the
full four-photoperiod calibration, as do all sourced parameters. The design
measures how much of the photoperiod response is carried by the growth
coefficient specifically, with the carbohydrate coefficients held fixed. It does
not establish that this coefficient is the only one that could carry that
response, and a design refitting all four per fold could distribute the same
dependence differently. A grid search is used rather than a local optimiser
because the objective is one-dimensional and bounded, so the grid resolution is
the only accuracy limit and no starting point or convergence criterion enters
the result.

The 18~h condition is not part of any fold. It is held out of every fit in
every regime, and it is evaluated only after the coefficients have been frozen.
It therefore tests extrapolation beyond the calibrated range, which is a
different and stronger question than the interpolation the four folds test.

\subsection{Ablation protocol}

The ablation measures what the computed carbon input contributes, by replacing
the carbon-input pathway and leaving everything downstream untouched. Three
variants are compared. In M0 the carbon input is prescribed by the simple
saturating photosynthesis expression of the reference formulation. In M1 the
carbon input comes from Farquhar C3 photosynthesis, with the supply--demand
closure solved as in the full model, but with stomatal conductance held fixed.
In M2 conductance is computed from the Jarvis multiplicative response, so the
gas-exchange layer becomes sensitive to PAR, vapour-pressure deficit,
temperature, ambient CO$_2$ and leaf water potential.

All parameters are shared and frozen across the three variants. This is the
point of the design and not a simplification. If each variant were re-fitted
against the same observations, the fit would absorb the difference between the
carbon-input pathways. The comparison would then measure how well a per-variant
calibration can compensate, rather than what removing a mechanism costs. The comparison is reported on RGR error per photoperiod. RGR alone cannot
separate M1 from M2, so we also report what the coupling changes. Two quantities
carry that: the magnitude and diel shape of the carbon input, and the response of
the gas-exchange layer to the environmental drivers.

\subsection{Sensitivity analysis protocol}

The global sensitivity analysis covers the 16 parameters listed in
\path{config/sensitivity_params.csv}. The admissible range, the sampling scale
and the literature basis of each are given individually in
Appendix~\ref{app:sa}. The set spans nutrient uptake capacity
and soil supply, gas-exchange capacity, the environmental drivers, respiration,
the growth coefficient, vacuolar exchange and the initial carbohydrate and
biomass state. The analysed output is RGR. Every reported index is conditional
on these declared ranges: the analysis describes model behaviour over the
declared domain and not parameter importance in general.

Morris screening follows \citet{Morris1991} with the elementary-effect measures
$\mu^\ast$ and $\sigma$ of \citet{Campolongo2007}. The design uses 30
trajectories on a 6-level grid with the fixed step $\Delta = p/[2(p-1)]$, giving
$r(d+1)$ model evaluations. Trajectory base points are drawn from an
independent pseudo-random generator, as the method specifies, so the screening
is not affected by the base-sample construction described below. Variance-based
indices use the estimators of \citet{Sobol2001} in the form recommended by
\citet{Saltelli2010}, with $n=512$ base samples and hence $n(d+2)$ evaluations
over the $d=16$ parameters. Confidence intervals come from 1000 bootstrap
replicates of the model evaluations, resampled by base row so that each
resampled design remains a valid Saltelli design.

The base-sample design is stated explicitly, because the estimator depends on
it. The Saltelli estimator needs two independent base matrices $\mathbf{A}$ and
$\mathbf{B}$ of size $n \times d$, and it attributes variance to a parameter
through a single-column substitution between them. Any correlation between the
corresponding columns therefore biases the index towards zero. The construction
specified by \citet{Saltelli2010} draws one low-discrepancy sequence in $2d$
dimensions and splits it by column, but the sampler used here supports at most
16 dimensions and $2d=32$ at $d=16$. The two matrices are therefore drawn as two
independent Latin-hypercube samples, stratified per dimension and independent by
construction \citep{McKay1979}. The design actually used is written to a
\texttt{base\_design} column of every output file, so no reported index is
ambiguous about its own provenance. The estimator implementation was checked
against the Ishigami function, whose indices are known in closed form
\citep{Ishigami1990}.

Uncertainty is reported at two levels, because the two levels answer different
questions. Bootstrap intervals resample within one design, so they measure
sampling noise given that design and are blind to error contributed by the
design itself. Across-design spread was therefore checked separately, by
repeating the whole analysis on independently generated base samples over base
sample sizes from $n=512$ to $n=4096$. An estimator whose across-design spread
does not shrink as $n$ grows is misspecified rather than noisy, so both levels
are needed before an index is interpreted. Divergent evaluations are counted and
reported with the indices, and no evaluation is silently discarded, because
deleting rows from a Saltelli design breaks the pairing the estimator relies on.
One-at-a-time sweeps and two-parameter response surfaces are also computed, but
they are local diagnostics of the response shape and are not used as a global
ranking.
